\section{Discusión y problemas abiertos}\label{sec:discusion}

\subsection{Una ilustración numérica}

En las seis filas de tamaño finito de la Tabla~\ref{tab:ilustracion} figuran los conteos enteros
$N_0$, $N_1$ y $N_2$, obtenidos por enumeración completa, y los valores de $f_1$ y $f_2$ que
resultan de ellos, redondeados a seis decimales. Los valores de $f_1(n)$ no son monótonos en
$n$. En la última fila, la existencia del límite común es el Teorema~\ref{thm:limite}, con
$F(0)=1/b_0$; el decimal que figura a su derecha es una estimación numérica de $1/b_0$, sin
intervalo de error certificado.

La estimación $F'(0)\approx0{,}516$ también es numérica. Sugiere comparar, a tamaño
finito, $f_h(n)$ con $F((h-\tfrac12)/n)$. La utilidad de esta comparación como
aproximación a tamaño finito es heurística. El resultado demostrado en el
Corolario~\ref{cor:capas} concierne a la suma en $K$ capas fijas cuando $n\to\infty$.
Los cocientes de la tabla se obtienen de los conteos mediante la ecuación~\eqref{eq:fh}.
Los conteos y el programa que los reproduce están en~\cite{Codigo}.

\begin{table}[ht]
\centering
\small
\begin{tabular}{rrrrccc}
\toprule
$n$ & $N_0(n)=Q(n)$ & $N_1(n)$ & $N_2(n)$ & $f_1(n)$ & $f_2(n)$ & $F(0)$ (numérico)\\
\midrule
10 & 724 & 56\,832 & 636\,524 & 0{,}127393 & 0{,}200891 & \\
11 & 2\,680 & 264\,544 & 3\,634\,976 & 0{,}111437 & 0{,}181943 & \\
12 & 14\,200 & 1\,375\,368 & 22\,146\,756 & 0{,}123894 & 0{,}170782 & \\
13 & 73\,712 & 7\,828\,784 & 141\,952\,580 & 0{,}122402 & 0{,}165452 & \\
14 & 365\,596 & 46\,788\,624 & 960\,059\,292 & 0{,}109393 & 0{,}158389 & \\
15 & 2\,279\,184 & 299\,307\,416 & 6\,849\,285\,616 & 0{,}114223 & 0{,}152947 & \\
$\infty$ & \multicolumn{5}{c}{límite común de $f_1$ y $f_2$ (Teorema~\ref{thm:limite})} & 0{,}080914\\
\bottomrule
\end{tabular}
\caption{Conteos exactos $N_0(n)=Q(n)$, $N_1(n)$ y $N_2(n)$, obtenidos por enumeración
exhaustiva; valores de $f_1$ y $f_2$, calculados con~\eqref{eq:fh} y redondeados a seis
decimales; y valor numérico del límite común. La columna $N_0$ coincide con la sucesión A000170
de~\cite{OEIS}, y $N_1(10)$ y $N_2(10)$ figuran también en el triángulo A348129 de esa
enciclopedia, que cuenta las colocaciones de $k$ reinas sin ataques en un tablero de $n\times n$.
El valor $0{,}080914$ es la
extrapolación de Richardson de segundo orden de los valores en las mallas de $128$ y $256$ puntos
($0{,}0809029$ y $0{,}0809114$); con las mallas de $64$ y $128$ puntos la extrapolación cambia en
menos de $10^{-8}$ (véase~\cite[Sección~9]{Perfil}). Es un valor numérico, sin certificación por
intervalos.}
\label{tab:ilustracion}
\end{table}

\subsection{Lo que queda abierto}

\begin{enumerate}[1.]
\item \emph{Tasa.} No se sabe si $f_h(n)-F(0)$, o $\log[f_h(n)/F((h-\frac12)/n)]$, es de orden
$1/n$. Las transferencias de la ley de referencia a la ley uniforme usan la cota de
entropía relativa $d_n=o(n)$ del Teorema~\ref{thm:comparacion}, sin tasa, y tolerancias fijas, de
modo que no pueden invertirse los límites para deducir una tasa.
\item \emph{Capas crecientes.} El término constante de $\log Q(n)$ depende de todas las capas.
Como $F$ se extiende a una función analítica y positiva en $[0,1]$~\cite[Teorema~8.1]{Perfil},
podemos considerar $\tau_{n,h}=n\log[f_h(n)/F((h-\frac12)/n)]$ para todas las capas,
$1\le h\le n$. Una ruta suficiente requiere la convergencia de los promedios
$\frac1n\sum_h\tau_{n,h}$ sobre los $h$ con $h/n$ en cada compacto de $(0,1)$ y un control
uniforme de los dos bordes; en el del tablero lleno, por ejemplo,
\[
 \lim_{x_0\downarrow0}\ \limsup_{n\to\infty}\ \sum_{h\le x_0n}
 \Bigl|\log\frac{f_h(n)}{F((h-\frac12)/n)}\Bigr|=0 .
\]
La convergencia para cada $h$ fijo no implica esta condición.
\item \emph{Ley conjunta de los huecos.} Para $h>1$ se identifica sólo la medida media de los
cruces vacíos. ¿Son asintóticamente independientes las posiciones de las filas y las columnas vacías? ¿Se
concentra
$d_{h,1}/n$ sin peso?
\item \emph{Constantes efectivas.} La prueba no da un umbral $n_0$ explícito ni constantes útiles
para tamaños moderados.
\end{enumerate}

Volviendo al tablero con $n-1$ reinas, la probabilidad de que el único cruce vacío
admita la última reina tiende a $1/b_0$. Si faltan $h$ reinas, con $h$ fijo, ese
mismo valor describe la probabilidad de añadir una reina en un cruce elegido al azar.
El problema que queda al dejar crecer $h$ es controlar cómo se acumulan las
diferencias entre la fracción $f_h(n)$ de cada capa y el valor correspondiente de $F$.
